\documentclass[12pt]{article}
\usepackage{fullpage}
\usepackage{amsmath,amssymb,amsthm}

\newtheorem{theorem}{Theorem}
\newtheorem{lemma}{Lemma}
\newtheorem{proposition}{Proposition}

\DeclareMathOperator{\Lk}{Lk}
\DeclareMathOperator{\St}{St}
\newcommand{\bbarotimes}{\mathbin{\overline{\otimes}}}

\begin{document}

\title{Proper proximality of irreducible graph-product von Neumann algebras}
\author{}
\date{}
\maketitle

\section*{Problem statement and interpretation}

I interpret ``irreducible'' in the standard graph-product sense: the finite
simplicial graph $\Gamma=(V,E)$ is not a non-trivial join, equivalently the
complement graph $\Gamma^c$ is connected.  Thus there is no partition
$V=S\sqcup T$ with $S,T\ne\varnothing$ such that every vertex of $S$ is adjacent
in $\Gamma$ to every vertex of $T$.  I also use the usual convention that the
traces $\varphi_v$ are faithful normal tracial states.

For $U\subset V$, let $M_U$ denote the graph product over the induced subgraph
$\Gamma|_U$; set $M_\varnothing=\mathbb C$.  We prove that
$M=M_V$ is properly proximal.

\section*{Standard inputs}

We use three standard facts.  The first two are the free-product and
amalgamated-free-product permanence results for proper proximality proved by
Ding--Kunnawalkam Elayavalli--Peterson using the Bass--Serre deformation and
the normal-bidual characterization of proper proximality \cite{DKEP}.  The
third is the relative amenability classification for graph-product von Neumann
algebras of Charlesworth et al.\ \cite{CDHJKN}.

\begin{theorem}[DKEP permanence]\label{thm:perm}
Let $B\subset P$ be an inclusion of separable tracial von Neumann algebras with
trace-preserving expectation, and let $(A,\tau_A)$ be a tracial von Neumann
algebra containing a unitary $u$ with $\tau_A(u)=0$.
\begin{enumerate}
\item If $P$ is properly proximal and $P$ is not amenable relative to $B$ inside
$P$, then
\[
        P *_B (B\bbarotimes A)
\]
is properly proximal.
\item If $A_0,A_1$ are tracial von Neumann algebras with trace-zero unitaries
and the pair is not the exceptional pair $(\mathbb C^2,\mathbb C^2)$ with
equal traces, then $A_0*A_1$ is properly proximal.
\end{enumerate}
\end{theorem}

\begin{theorem}[Relative amenability in graph products]\label{thm:relam}
Let $X$ be a finite vertex set and let
$N=M_X$ be the graph product of $\{M_x:x\in X\}$, each $M_x$ containing a
trace-zero unitary.  For $Y\subset X$, $N$ is amenable relative to $M_Y$ inside
$N$ if and only if the following two conditions hold:
\begin{enumerate}
\item $M_x$ is amenable for every $x\in X\setminus Y$;
\item whenever $x\in X\setminus Y$ and $z\in X\setminus\{x\}$, either
$x$ and $z$ are adjacent in $\Gamma|_X$, or else
\[
\dim M_x=\dim M_z=2
\]
and both $x$ and $z$ are adjacent to every vertex of
$X\setminus\{x,z\}$.
\end{enumerate}
\end{theorem}

We also use the standard graph-product decomposition: if $W\subset V$,
$v\notin W$, and $L=W\cap \Lk_\Gamma(v)$, then
\begin{equation}\label{eq:afp}
        M_{W\cup\{v\}}
        \cong
        M_W *_{M_L} \bigl(M_L\bbarotimes M_v\bigr).
\end{equation}
Indeed, $M_v$ commutes precisely with $M_L$ among the old vertex algebras, and
the reduced-word description of graph products gives freeness over $M_L$.

\section*{Proof}

Since $\Gamma$ is irreducible, $\Gamma^c$ is connected.  Choose a path
$v_1-v_2-v_3$ in $\Gamma^c$.  Extend it to an ordering
\[
        V=\{v_1,\ldots,v_n\}
\]
such that, for each $k\ge4$, the vertex $v_k$ has a neighbor
$r_k\in \{v_1,\ldots,v_{k-1}\}$ in $\Gamma^c$, and $r_k$ has a neighbor
$s_k\in\{v_1,\ldots,v_{k-1}\}$ in $\Gamma^c$.  This is obtained, for example,
from a spanning tree of $\Gamma^c$ rooted at $v_2$ and then listing vertices by
nondecreasing distance from the root.  In particular, the induced complement
graph on $V_k:=\{v_1,\ldots,v_k\}$ is connected for every $k\ge3$.

We first prove proper proximality for $M_{V_3}$.  In $\Gamma$, the vertex
$v_2$ is not adjacent to $v_1$ or $v_3$, hence
\[
        M_{V_3}\cong M_{\{v_1,v_3\}} * M_{v_2}.
\]
The algebra $M_{\{v_1,v_3\}}$ contains a trace-zero unitary and is not
two-dimensional: if $v_1$ and $v_3$ are adjacent in $\Gamma$ it is
$M_{v_1}\bbarotimes M_{v_3}$, while otherwise it is the free product
$M_{v_1}*M_{v_3}$; in either case both vertex algebras are non-scalar.  Thus
Theorem \ref{thm:perm}(2) gives that $M_{V_3}$ is properly proximal.

Assume inductively that $P:=M_{V_{k-1}}$ is properly proximal for some
$k\ge4$.  Put
\[
        L=V_{k-1}\cap \Lk_\Gamma(v_k).
\]
By \eqref{eq:afp},
\[
        M_{V_k}\cong P*_{M_L}(M_L\bbarotimes M_{v_k}).
\]
The new vertex algebra $M_{v_k}$ has a trace-zero unitary, so Theorem
\ref{thm:perm}(1) will apply once we know that $P$ is not amenable relative to
$M_L$.

Suppose, toward a contradiction, that $P=M_{V_{k-1}}$ is amenable relative to
$M_L$.  Since $r_k$ is adjacent to $v_k$ in $\Gamma^c$, it is not in $L$.
Since $r_k$ is adjacent to $s_k$ in $\Gamma^c$, the vertices $r_k$ and $s_k$
are not adjacent in $\Gamma$.  Applying Theorem \ref{thm:relam} to
$X=V_{k-1}$ and $Y=L$, condition (2) forces the exceptional alternative:
$r_k$ and $s_k$ must both be adjacent in $\Gamma$ to every vertex of
$V_{k-1}\setminus\{r_k,s_k\}$.  Equivalently, in the complement graph induced
on $V_{k-1}$, the only complement-neighbor of $r_k$ or $s_k$ among
$V_{k-1}$ is the other one.  Thus $\{r_k,s_k\}$ is a connected component of
$(\Gamma|_{V_{k-1}})^c$.  This contradicts the connectedness of
$(\Gamma|_{V_{k-1}})^c$ and the fact that $|V_{k-1}|\ge3$.

Therefore $P$ is not amenable relative to $M_L$, and Theorem
\ref{thm:perm}(1) implies that $M_{V_k}$ is properly proximal.  Induction gives
proper proximality of $M_{V_n}=M$, as required.

\begin{thebibliography}{9}

\bibitem{DKEP}
C.~Ding, S.~Kunnawalkam Elayavalli, and J.~Peterson,
\emph{Properly proximal von Neumann algebras},
Duke Math. J. \textbf{172} (2023), 2821--2894.

\bibitem{CDHJKN}
I.~Charlesworth, R.~de Santiago, B.~Hayes, D.~Jekel,
S.~Kunnawalkam Elayavalli, and B.~Nelson,
\emph{On the structure of graph product von Neumann algebras},
Publ. Res. Inst. Math. Sci. \textbf{61} (2025), 713--762.

\end{thebibliography}

\end{document}
